\documentclass[11pt]{article}
\usepackage[utf8]{inputenc}
\usepackage{amsmath,amssymb}
\usepackage[margin=1in]{geometry}
\usepackage{hyperref}
\emergencystretch=3em

\title{The zero-phase dictionary and the constant cross-validation}
\author{Claude, with Daniel Murfet}
\date{2026-09-07}

\begin{document}
\maketitle
\sloppy

\begin{abstract}
Astra~\#18's ``dictionary'': the two-dimensional quadratic chart of the random-phase Taylor tree,
at zero phase $\xi\equiv0$ and unit amplitude $\eta\equiv1$ (Taylor data $x=0$, $y=\delta_{00}$),
is \emph{exactly} the unit-box monomial integral at parameter $N^2$
(\texttt{twoDAmp\_zeroPhase\_eq\_monomial\_sq}). A square-parameter transport lemma turns a
power--log limit in $T$ into one in $N=\sqrt T$ with the logarithmic coefficient multiplied by $2^r$
(\texttt{tendsto\_powLog\_comp\_sq}). By uniqueness of limits the Taylor-tree leading log coefficient
at zero phase equals \emph{twice} the monomial mixed constant at $\lambda=p/2$
(\texttt{canonA\_zeroPhase\_eq\_two\_mul\_monomialMixedConst}):
\[
\frac{y_{00}}{k_1k_2}\int_0^\infty s^{p-1}e^{-\beta s^2}\,ds
=\frac{y_{00}\,\Gamma(p/2)\beta^{-p/2}}{2k_1k_2}
=2\cdot\frac{\Gamma(p/2)\beta^{-p/2}}{1!\cdot2k_1\cdot2k_2},
\]
the factor $2$ being $\log N^2=2\log N$. This cross-validates the constants of the two programmes at
the level of theorems. Zero \texttt{sorry}/\texttt{axiom}.
\end{abstract}

\section{The things formalised}
{\small
\begin{verbatim}
theorem ampCoeff_zero_delta (β) (k) (s) : ampCoeff β (fun _ => 0) delta k s = delta k
theorem anaAmp_ampCoeff_zero_delta (β b u v s) :
    anaAmp (ampCoeff β (fun _ => 0) delta) b u v s = 1
theorem twoDAmp_const_eq_monomialCutoff (β b N) (h₁ h₂ k₁ k₂) :
    twoDAmp β b N h₁ h₂ k₁ k₂ (fun _ _ _ => 1)
      = monomialBoxRealCutoff 2 ![h₁, h₂] ![k₁, k₂] b β (N ^ 2)
theorem twoDAmp_zeroPhase_eq_monomial_sq (β N) (h₁ h₂ k₁ k₂) :
    twoDAmp β 1 N h₁ h₂ k₁ k₂ (anaAmp (ampCoeff β (fun _ => 0) delta) 1)
      = monomialBoxReal 2 ![h₁, h₂] ![k₁, k₂] β (N ^ 2)
theorem tendsto_powLog_comp_sq (F) (l C) (r)
    (hF : Tendsto (fun T => F T / (T ^ (-l) * log T ^ r)) atTop (𝓝 C)) :
    Tendsto (fun N => F (N ^ 2) / (N ^ (-(2 * l)) * log N ^ r)) atTop (𝓝 (2 ^ r * C))
theorem canonA_zeroPhase_eq_two_mul_monomialMixedConst (β) (h₁ h₂ k₁ k₂) (hβ) (hk₁) (hk₂) (p)
    (hp₁ : ((h₁ : ℝ) + 1) / k₁ = p) (hp₂ : ((h₂ : ℝ) + 1) / k₂ = p) :
    canonA β h₁ h₂ k₁ k₂ (fun i j s => ampCoeff β (fun _ => 0) delta (i, j) s) p
      = 2 * monomialMixedConst ![h₁, h₂] ![k₁, k₂] (p / 2) β
\end{verbatim}
}

\section{Mathematics}
The exponential coefficient array of the zero phase is $\delta_{00}$ (only the $n=0$ term of the
exponential series survives), so the amplitude array is $\delta_{00}$ and its analytic amplitude is
the constant $1$; the quadratic kernel $e^{-\beta(Nu^{k_1}v^{k_2})^2}$ is $e^{-\beta N^2u^{2k_1}v^{2k_2}}$,
and two Bochner peels identify the iterated integral with the $\mathrm{Fin}\,2$ box integral. The
Taylor tree gives $\mathrm{twoDAmp}\sim A\,N^{-p}\log N$ with $A>0$ (units 154, 167); the monomial
theorem at $T=N^2$ gives the same ratio with limit $2\cdot\Gamma(p/2)\beta^{-p/2}/(2k_1\cdot2k_2)$.
Consistency with the fluctuation-function normalisation: $J_p(0)=\tfrac12\beta^{-p/2}\Gamma(p/2)$
(\texttt{gaussMomentJ\_zero}), $S_{p/2}(a)=2J_p(a)$, $A_p=y_{00}J_p(x_{00})/(k_1k_2)$.

\section{Paper-facing meaning}
The two independent developments --- the random-phase Taylor tree in the variable $N=\sqrt n$ and
the monomial normal-moment programme in the variable $N=n$ --- agree on the leading constant where
they overlap. In the paper's variable $n$: $Q(\sqrt n)/(n^{-p/2}\log n)\to y_{00}\Gamma(p/2)\beta^{-p/2}/(4k_1k_2)$,
i.e.\ the coefficient of $N^{-p}\log N$ is twice the coefficient of $n^{-p/2}\log n$.

\section{Lean notes}
\begin{itemize}
\item \texttt{isEquivalent\_iff\_tendsto\_one} now takes \texttt{∀ᶠ x, v x ≠ 0} (not the older
implication form).
\item Indices of \texttt{Fin 2}: \texttt{Fin.cons u (Fin.cons v \_) 1} needs
\texttt{← Fin.succ\_zero\_eq\_one, Fin.cons\_succ, Fin.cons\_zero}; \texttt{Matrix.cons\_val\_one}
handles \texttt{![h₁, h₂] 1}.
\item Continuity of \texttt{fun w => Fin.cons u w} by \texttt{continuous\_pi} and \texttt{Fin.cases}.
\item A \texttt{simp} on a sum of indicator-like \texttt{if}s over \texttt{Fin 2} may rewrite it into
a \texttt{Finset.card} of a set-builder; use \texttt{Fin.sum\_univ\_two} and \texttt{if\_pos} instead.
\end{itemize}

\end{document}
